\chapter{Primary Decomposition}\label{chapter4}

The decomposition of an ideal into primary ideals is a traditional pillar of ideal theory. It provides the algebraic foundation for decomposing an algebraic variety into its irreducible components---although it is only fair to point out that the algebraic picture is more complicated than naive geometry would suggest. From another point of view primary decomposition provides a generalization of the factorization of an integer as a product of prime-powers. In the modern treatment, with its emphasis on localization, primary decomposition is no longer such a central tool in the theory. It is still, however, of interest in itself and in this chapter we establish the classical uniqueness theorems.

The prototypes of commutative rings are $\Z$ and the ring of polynomials $k[x_1, \ldots, x_n]$ where $k$ is a field; both these are unique factorization domains. This is not true of arbitrary commutative rings, even if they are integral domains (the classical example is the ring $\Z[\sqrt{-5}]$, in which the element $6$ has two essentially distinct factorizations, $2 \cdot 3$ and $(1 + \sqrt{-5})(1 - \sqrt{-5})$). However, there is a generalized form of ``unique factorization'' of ideals (not of elements) in a wide class of rings (the Noetherian rings).

A prime ideal in a ring $A$ is in some sense a generalization of a prime number. The corresponding generalization of a power of a prime number is a primary ideal. An ideal $\mathfrak{q}$ in a ring $A$ is primary if $\mathfrak{q} \neq A$ and if
\[
    xy \in \mathfrak{q} \implies \text{either } x \in \mathfrak{q} \text{ or } y^n \in \mathfrak{q} \text{ for some } n > 0.
\]
In other words,
\[
    \mathfrak{q} \text{ is primary} \iff A/\mathfrak{q} \neq 0 \text{ and every zero-divisor in } A/\mathfrak{q} \text{ is nilpotent.}
\]

Clearly every prime ideal is primary. Also the contraction of a primary ideal is primary, for if $f \colon A \to B$ and if $\mathfrak{q}$ is a primary ideal in $B$, then $A/\mathfrak{q}^c$ is isomorphic to a subring of $B/\mathfrak{q}$.

\begin{proposition}\label{prop:radical-primary}
Let $\mathfrak{q}$ be a primary ideal in a ring $A$. Then $r(\mathfrak{q})$ is the smallest prime ideal containing $\mathfrak{q}$.
\end{proposition}

\begin{proof}
    By \eqref{prop:nilradical-intersection} it is enough to show that $\mathfrak{p} = r(\mathfrak{q})$ is prime. Let $xy \in r(\mathfrak{q})$, then $(xy)^m \in \mathfrak{q}$ for some $m > 0$, and therefore either $x^m \in \mathfrak{q}$ or $y^{mn} \in \mathfrak{q}$ for some $n > 0$; i.e., either $x \in r(\mathfrak{q})$ or $y \in r(\mathfrak{q})$.
\end{proof}

If $\mathfrak{p} = r(\mathfrak{q})$, then $\mathfrak{q}$ is said to be \emph{$\mathfrak{p}$-primary}\index{p-primary ideal@$\mathfrak{p}$-primary ideal}.

\begin{examples}
    \begin{enumerate}[(1)]
        \item The primary ideals in $\Z$ are $(0)$ and $(p^n)$, where $p$ is prime. For these are the only ideals in $\Z$ with prime radical, and it is immediately checked that they are primary.
        
        \item Let $A = k[x, y]$, $\mathfrak{q} = (x, y^2)$. Then $A/\mathfrak{q} \cong k[y]/(y^2)$, in which the zero-divisors are all the multiples of $y$, hence are nilpotent. Hence $\mathfrak{q}$ is primary, and its radical $\mathfrak{p}$ is $(x, y)$. We have $\mathfrak{p}^2 \subset \mathfrak{q} \subset \mathfrak{p}$ (strict inclusions), so that a primary ideal is not necessarily a prime-power.
        
        \item Conversely, a prime power $\mathfrak{p}^n$ is not necessarily primary, although its radical is the prime ideal $\mathfrak{p}$. For example, let $A = k[x, y, z]/(xy - z^2)$ and let $\bar{x}$, $\bar{y}$, $\bar{z}$ denote the images of $x$, $y$, $z$ respectively in $A$. Then $\mathfrak{p} = (\bar{x}, \bar{z})$ is prime (since $A/\mathfrak{p} \cong k[y]$, an integral domain); we have $\bar{x}\bar{y} = \bar{z}^2 \in \mathfrak{p}^2$ but $\bar{x} \notin \mathfrak{p}^2$ and $\bar{y} \notin r(\mathfrak{p}^2) = \mathfrak{p}$; hence $\mathfrak{p}^2$ is not primary.
    \end{enumerate}
\end{examples}

However, there is the following result:

\begin{proposition}\label{prop:radical-maximal}
    If $r(\mathfrak{a})$ is maximal, then $\mathfrak{a}$ is primary. In particular, the powers of a maximal ideal $\mathfrak{m}$ are $\mathfrak{m}$-primary.
\end{proposition}

\begin{proof}
    Let $r(\mathfrak{a}) = \mathfrak{m}$. The image of $\mathfrak{m}$ in $A/\mathfrak{a}$ is the nilradical of $A/\mathfrak{a}$, hence $A/\mathfrak{a}$ has only one prime ideal, by \eqref{prop:nilradical-intersection}. Hence every element of $A/\mathfrak{a}$ is either a unit or nilpotent, and so every zero-divisor in $A/\mathfrak{a}$ is nilpotent.
\end{proof}

We are going to study presentations of an ideal as an intersection of primary ideals. First, a couple of lemmas:

\begin{lemma}\label{lem:intersection-primary}
    If $\mathfrak{q}_i$ ($1 \leq i \leq n$) are $\mathfrak{p}$-primary, then $\mathfrak{q} = \bigcap_{i=1}^n \mathfrak{q}_i$ is $\mathfrak{p}$-primary.
\end{lemma}

\begin{proof}
    $r(\mathfrak{q}) = r(\bigcap_{i=1}^n \mathfrak{q}_i) = \bigcap r(\mathfrak{q}_i) = \mathfrak{p}$. Let $xy \in \mathfrak{q}$, $y \notin \mathfrak{q}$. Then for some $i$ we have $xy \in \mathfrak{q}_i$ and $y \notin \mathfrak{q}_i$, hence $x \in \mathfrak{p}$, since $\mathfrak{q}_i$ is primary.
\end{proof}

\begin{lemma}\label{lem:colon-primary}\index{colon ideal}
Let $\mathfrak{q}$ be a $\mathfrak{p}$-primary ideal, $x$ an element of $A$. Then
\begin{enumerate}[i)]
    \item if $x \in \mathfrak{q}$ then $(\mathfrak{q} : x) = (1)$;
    \item if $x \notin \mathfrak{q}$ then $(\mathfrak{q} : x)$ is $\mathfrak{p}$-primary, and therefore $r(\mathfrak{q} : x) = \mathfrak{p}$;
    \item if $x \notin \mathfrak{p}$ then $(\mathfrak{q} : x) = \mathfrak{q}$.
\end{enumerate}
\end{lemma}

\begin{proof}
(i) and (iii) follow immediately from the definitions.

(ii): if $y \in (\mathfrak{q} : x)$ then $xy \in \mathfrak{q}$, hence (as $x \notin \mathfrak{q}$) we have $y \in \mathfrak{p}$. Hence $\mathfrak{q} \subseteq (\mathfrak{q} : x) \subseteq \mathfrak{p}$; taking radicals, we get $r(\mathfrak{q} : x) = \mathfrak{p}$. Let $yz \in (\mathfrak{q} : x)$ with $y \notin \mathfrak{p}$; then $xyz \in \mathfrak{q}$, hence $xz \in \mathfrak{q}$, hence $z \in (\mathfrak{q} : x)$.
\end{proof}

A \emph{primary decomposition}\index{primary decomposition} of an ideal $\mathfrak{a}$ in $A$ is an expression of $\mathfrak{a}$ as a finite intersection of primary ideals, say
\begin{equation}\label{eq:primary-decomposition}
    \mathfrak{a} = \bigcap_{i=1}^n \mathfrak{q}_i. \tag{1}
\end{equation}
(In general such a primary decomposition need not exist; in this chapter we shall restrict our attention to ideals which have a primary decomposition.) If moreover (i) the $r(\mathfrak{q}_i)$ are all distinct, and (ii) we have $\mathfrak{q}_i \not\supseteq \bigcap_{j \neq i} \mathfrak{q}_j$ ($1 \leq i \leq n$) the primary decomposition~\eqref{eq:primary-decomposition} is said to be \emph{minimal} (or \emph{irredundant}, or \emph{reduced}, or \emph{normal}, \ldots). By \eqref{lem:intersection-primary} we can achieve (i) and then we can omit any superfluous terms to achieve (ii); thus any primary decomposition can be reduced to a minimal one. We shall say that $\mathfrak{a}$ is \emph{decomposable}\index{decomposable ideal} if it has a primary decomposition.

\begin{theorem}[1st Uniqueness Theorem]\label{thm:first-uniqueness}\index{First Uniqueness Theorem}\index{associated prime ideals}
    Let $\mathfrak{a}$ be a decomposable ideal and let $\mathfrak{a} = \bigcap_{i=1}^n \mathfrak{q}_i$ be a minimal primary decomposition of $\mathfrak{a}$. Let $\mathfrak{p}_i = r(\mathfrak{q}_i)$ $(1 \leq i \leq n)$. Then the $\mathfrak{p}_i$ are precisely the prime ideals which occur in the set of ideals $r(\mathfrak{a} : x)$ $(x \in A)$, and hence are independent of the particular decomposition of $\mathfrak{a}$.
\end{theorem}

\begin{proof}
    For any $x \in A$ we have $(\mathfrak{a} : x) = \left(\bigcap \mathfrak{q}_i : x\right) = \bigcap (\mathfrak{q}_i : x)$, hence $r(\mathfrak{a} : x) = \bigcap_{i=1}^n r(\mathfrak{q}_i : x) = \bigcap_{x \notin \mathfrak{q}_i} \mathfrak{p}_i$ by \eqref{lem:colon-primary}. Suppose $r(\mathfrak{a} : x)$ is prime; then by \eqref{pro:prime avoidance} we have $r(\mathfrak{a} : x) = \mathfrak{p}_j$ for some $j$. Hence every prime ideal of the form $r(\mathfrak{a} : x)$ is one of the $\mathfrak{p}_i$. Conversely, for each $i$ there exists $x_i \notin \mathfrak{q}_i$, $x_i \in \bigcap_{j \neq i} \mathfrak{q}_j$, since the decomposition is minimal; and we have $r(\mathfrak{a} : x_i) = \mathfrak{p}_i$.
\end{proof}
\begin{remarks}\label{rem:first-uniqueness}
    \begin{enumerate}[1)]
        \item The above proof, coupled with the last part of \eqref{lem:colon-primary}, shows that for each $i$ there exists $x_i$ in $A$ such that $(\mathfrak{a} : x_i)$ is $\mathfrak{p}_i$-primary.
        \item Considering $A/\mathfrak{a}$ as an $A$-module, \eqref{thm:first-uniqueness} is equivalent to saying that the $\mathfrak{p}_i$ are precisely the prime ideals which occur as radicals of annihilators of elements of $A/\mathfrak{a}$.
    \end{enumerate}
\end{remarks}

\begin{example}\label{ex:embedded-primes}
    Let $\mathfrak{a} = (x^2, xy)$ in $A = k[x, y]$. Then $\mathfrak{a} = \mathfrak{p}_1 \cap \mathfrak{p}_2^2$ where $\mathfrak{p}_1 = (x)$, $\mathfrak{p}_2 = (x, y)$. The ideal $\mathfrak{p}_2^2$ is primary by \eqref{prop:radical-maximal}. So the prime ideals are $\mathfrak{p}_1$, $\mathfrak{p}_2$. In this example $\mathfrak{p}_1 \subset \mathfrak{p}_2$; we have $r(\mathfrak{a}) = \mathfrak{p}_1 \cap \mathfrak{p}_2 = \mathfrak{p}_1$, but $\mathfrak{a}$ is not a primary ideal.
\end{example}

The prime ideals $\mathfrak{p}_i$ in Theorem~\ref{thm:first-uniqueness} are said to \emph{belong} to $\mathfrak{a}$, or to be \emph{associated} with $\mathfrak{a}$. The ideal $\mathfrak{a}$ is primary if and only if it has only one associated prime ideal. The minimal elements of the set $\{\mathfrak{p}_1, \ldots, \mathfrak{p}_n\}$ are called the \emph{minimal} or \emph{isolated}\index{isolated prime ideal} prime ideals belonging to $\mathfrak{a}$. The others are called \emph{embedded}\index{embedded prime ideal} prime ideals. In the example above, $\mathfrak{p}_2 = (x, y)$ is embedded.

\begin{proposition}\label{prop:minimal-primes}
    Let $\mathfrak{a}$ be a decomposable ideal. Then any prime ideal $\mathfrak{p} \supseteq \mathfrak{a}$ contains a minimal prime ideal belonging to $\mathfrak{a}$, and thus the minimal prime ideals of $\mathfrak{a}$ are precisely the minimal elements in the set of all prime ideals containing $\mathfrak{a}$.
\end{proposition}

\begin{proof}
    If $\mathfrak{p} \supseteq \mathfrak{a} = \bigcap_{i=1}^n \mathfrak{q}_i$, then $\mathfrak{p} = r(\mathfrak{p}) \supseteq \bigcap r(\mathfrak{q}_i) = \bigcap \mathfrak{p}_i$. Hence by \eqref{pro:prime avoidance} we have $\mathfrak{p} \supseteq \mathfrak{p}_i$ for some $i$; hence $\mathfrak{p}$ contains a minimal prime ideal of $\mathfrak{a}$.
\end{proof}

\begin{remarks}
    \begin{enumerate}[1)]
        \item The names \emph{isolated} and \emph{embedded} come from geometry. Thus if $A = k[x_1, \ldots, x_n]$ where $k$ is a field, the ideal $\mathfrak{a}$ gives rise to a variety $X \subseteq k^n$ (see Chapter \ref{chapter1}, Exercise \ref{exer25-chapter1}). The minimal primes $\mathfrak{p}_i$ correspond to the irreducible components of $X$, and the embedded primes correspond to subvarieties of these, i.e., varieties \emph{embedded} in the irreducible components. Thus in the example before \eqref{prop:minimal-primes} the variety defined by $\mathfrak{a}$ is the line $x=0$, and the embedded ideal $\mathfrak{p}_2=(x,y)$ corresponds to the origin $(0,0)$.
        \item It is \emph{not} true that all the primary components are independent of the decomposition. For example $(x^2,xy)=(x)\cap (x,y)^2=(x)\cap (x^2,y)$ are two distinct primary decompositions. However, there are some uniqueness properties: see \eqref{thm:second-uniqueness}.
    \end{enumerate}
\end{remarks}

\begin{proposition}\label{pro:zero-divisors-union-primes}
    Let $\mathfrak{a}$ be a decomposable ideal, let $\mathfrak{a} = \bigcap_{i=1}^n \mathfrak{q}_i$ be a minimal primary decomposition, and let $r(\mathfrak{q}_i) = \mathfrak{p}_i$. Then
    \[
    \bigcup_{i=1}^n \mathfrak{p}_i = \{ x \in A : (\mathfrak{a} : x) \neq \mathfrak{a} \}.
    \]
    In particular, if the zero ideal is decomposable, the set $D$ of zero-divisors of $A$ is the union of the prime ideals belonging to $0$.
\end{proposition}

\begin{proof}
    If $\mathfrak{a}$ is decomposable, then $0$ is decomposable in $A/\mathfrak{a}$: namely $0 = \bigcap \overline{\mathfrak{q}}_i$ where $\overline{\mathfrak{q}}_i$ is the image of $\mathfrak{q}_i$ in $A/\mathfrak{a}$, and is primary. Hence it is enough to prove the last statement of \eqref{pro:zero-divisors-union-primes}. By \eqref{prop:zerodivisors-union} we have $D = \bigcup_{x \neq 0} r(0:x)$; from the proof of \eqref{thm:first-uniqueness}, we have $r(0:x) = \bigcap_{x \notin \mathfrak{q}_j} \mathfrak{p}_j \subseteq \mathfrak{p}_j$ for some $j$, hence $D \subseteq \bigcup_{i=1}^n \mathfrak{p}_i$. But also from \eqref{thm:first-uniqueness} each $\mathfrak{p}_i$ is of the form $r(0:x)$ for some $x \in A$, hence $\bigcup \mathfrak{p}_i \subseteq D$.
\end{proof}

Thus (the zero ideal being decomposable)
\[
\begin{aligned}
D &= \text{set of zero-divisors} \\
&= \bigcup \text{ of all prime ideals belonging to } 0; \\
\mathfrak{N} &= \text{set of nilpotent elements} \\
&= \bigcap \text{ of all minimal primes belonging to } 0.
\end{aligned}
\]

Next we investigate the behavior of primary ideals under localization.

\begin{proposition}\label{prop:localization-primary}\index{localization!of primary ideal}
    Let $S$ be a multiplicatively closed subset of $A$ and $\mathfrak{q}$ a $\mathfrak{p}$-primary ideal.
    \begin{enumerate}[i)]
        \item If $S \cap \mathfrak{p} \neq \varnothing$, then $S^{-1}\mathfrak{q} = S^{-1}A$.
        \item If $S \cap \mathfrak{p} = \varnothing$, then $S^{-1}\mathfrak{q}$ is $S^{-1}\mathfrak{p}$-primary in $S^{-1}A$, and its contraction to $A$ is $\mathfrak{q}$.
    \end{enumerate}

    Hence primary ideals correspond to primary ideals in the correspondence \eqref{prop:extension-ideals} between ideals in $S^{-1}A$ and contracted ideals in $A$.
\end{proposition}
\begin{proof}
    \begin{enumerate}[i)]
        \item If $s \in S \cap \mathfrak{p}$, then $s^n \in \mathfrak{q}$ for some $n$, hence $s^n/1 \in S^{-1}\mathfrak{q}$. Since $s/1$ is a unit in $S^{-1}A$, it follows that $S^{-1}\mathfrak{q} = S^{-1}A$.
        \item If $S \cap \mathfrak{p} = \varnothing$, then $s \in S$ and $as \in \mathfrak{q}$ imply $a \in \mathfrak{q}$, hence $\mathfrak{q}^{ec} = \mathfrak{q}$ by \eqref{prop:extension-ideals}.
        Also from \eqref{prop:extension-ideals} we have $r(\mathfrak{q}^e) = r(S^{-1}\mathfrak{q}) = S^{-1}r(\mathfrak{q}) = S^{-1}\mathfrak{p}$. The verification that $S^{-1}\mathfrak{q}$ is primary is straightforward. Finally, the contraction of a primary ideal is primary. 
    \end{enumerate}
\end{proof}

For any ideal $\mathfrak{a}$ and any multiplicatively closed subset $S$ in $A$, the contraction in $A$ of the ideal $S^{-1}\mathfrak{a}$ is denoted by $S(\mathfrak{a})$.

\begin{proposition}\label{prop:localization-primary-decomposition}
    Let $S$ be a multiplicatively closed subset of $A$ and let $\mathfrak{a}$ be a decomposable ideal. Let $\mathfrak{a} = \bigcap_{i=1}^n \mathfrak{q}_i$ be a minimal primary decomposition of $\mathfrak{a}$. Let $\mathfrak{p}_i = r(\mathfrak{q}_i)$ and suppose the $\mathfrak{q}_i$ numbered so that $S$ meets $\mathfrak{p}_{m+1}, \dots, \mathfrak{p}_n$ but not $\mathfrak{p}_1, \dots, \mathfrak{p}_m$. Then
    \[
    S^{-1}\mathfrak{a} = \bigcap_{i=1}^m S^{-1}\mathfrak{q}_i, \quad S(\mathfrak{a}) = \bigcap_{i=1}^m \mathfrak{q}_i,
    \]
    and these are minimal primary decompositions.
\end{proposition}

\begin{proof}
    $S^{-1}\mathfrak{a} = \bigcap_{i=1}^n S^{-1}\mathfrak{q}_i$ by \eqref{prop:extension-ideals} $= \bigcap_{i=1}^m S^{-1}\mathfrak{q}_i$ by \eqref{prop:localization-primary}, and $S^{-1}\mathfrak{q}_i$ is $S^{-1}\mathfrak{p}_i$-primary for $i = 1, \dots, m$. Since the $\mathfrak{p}_i$ are distinct, so are the $S^{-1}\mathfrak{p}_i$ ($1 \le i \le m$), hence we have a minimal primary decomposition. Contracting both sides, we get
    \[
    S(\mathfrak{a}) = (S^{-1}\mathfrak{a})^c = \bigcap_{i=1}^m (S^{-1}\mathfrak{q}_i)^c = \bigcap_{i=1}^m \mathfrak{q}_i
    \]
    by \eqref{prop:localization-primary} again. 
\end{proof}

A set $\Sigma$ of prime ideals belonging to $\mathfrak{a}$ is said to be \emph{isolated}\index{isolated prime ideal} if it satisfies the following condition: if $\mathfrak{p}'$ is a prime ideal belonging to $\mathfrak{a}$ and $\mathfrak{p}' \subseteq \mathfrak{p}$ for some $\mathfrak{p} \in \Sigma$, then $\mathfrak{p}' \in \Sigma$.

Let $\Sigma$ be an isolated set of prime ideals belonging to $\mathfrak{a}$, and let $S = A - \bigcup_{\mathfrak{p} \in \Sigma} \mathfrak{p}$. Then $S$ is multiplicatively closed and, for any prime ideal $\mathfrak{p}'$ belonging to $\mathfrak{a}$, we have
\begin{align*}
\mathfrak{p}' \in \Sigma &\implies \mathfrak{p}' \cap S = \varnothing; \\
\mathfrak{p}' \notin \Sigma &\implies \mathfrak{p}' \not\subseteq \bigcup_{\mathfrak{p} \in \Sigma} \mathfrak{p} \ (\text{by \eqref{pro:prime avoidance}}) \implies \mathfrak{p}' \cap S \neq \varnothing.
\end{align*}
Hence, from \eqref{prop:localization-primary-decomposition}, we deduce

\begin{theorem}[2nd uniqueness theorem]\label{thm:second-uniqueness}
    Let $\mathfrak{a}$ be a decomposable ideal, let $\mathfrak{a} = \bigcap_{i=1}^n \mathfrak{q}_i$ be a minimal primary decomposition of $\mathfrak{a}$, and let $\{\mathfrak{p}_{i_1}, \dots, \mathfrak{p}_{i_m}\}$ be an isolated set of prime ideals of $\mathfrak{a}$. Then $\mathfrak{q}_{i_1} \cap \cdots \cap \mathfrak{q}_{i_m}$ is independent of the decomposition.
\end{theorem}

In particular:

\begin{corollary}\label{cor:isolated-primary-components-unique}
    The isolated primary components (i.e., the primary components $\mathfrak{q}_i$ corresponding to minimal prime ideals $\mathfrak{p}_i$) are uniquely determined by $\mathfrak{a}$.
\end{corollary}

\begin{proof}[Proof of \eqref{thm:second-uniqueness}]
    We have $\mathfrak{q}_{i_1} \cap \cdots \cap \mathfrak{q}_{i_m} = S(\mathfrak{a})$ where $S = A - \mathfrak{p}_{i_1} \cup \cdots \cup \mathfrak{p}_{i_m}$, hence depends only on $\mathfrak{a}$ (since the $\mathfrak{p}_i$ depend only on $\mathfrak{a}$). 
\end{proof}

\begin{remark}
On the other hand, the embedded primary components are in general not uniquely determined by $\mathfrak{a}$. If $A$ is a Noetherian ring, there are in fact infinitely many choices for each embedded component (see Chapter \ref{chapter8}, Exercise \ref{exer1-chapter8}).
\end{remark}

%=============================================================================
% Chapter 4 Exercises - With hints from original book (23 exercises)
%=============================================================================
\section*{Exercises}
\addcontentsline{toc}{section}{Exercises}

\begin{enumerate}
    \item \label{exer1-chapter4}If an ideal $\mathfrak{a}$ has a primary decomposition, then $\Spec(A/\mathfrak{a})$ has only finitely many irreducible components.

    \item \label{exer2-chapter4}If $\mathfrak{a} = r(\mathfrak{a})$, then $\mathfrak{a}$ has no embedded prime ideals.

    \item \label{exer3-chapter4}If $A$ is absolutely flat, every primary ideal is maximal.

    \item \label{exer4-chapter4}In the polynomial ring $\Z[t]$, the ideal $\mathfrak{m} = (2, t)$ is maximal and the ideal $\mathfrak{q} = (4, t)$ is $\mathfrak{m}$-primary, but is not a power of $\mathfrak{m}$.

    \item \label{exer5-chapter4}In the polynomial ring $K[x, y, z]$ where $K$ is a field and $x, y, z$ are independent indeterminates, let $\mathfrak{p}_1 = (x, y)$, $\mathfrak{p}_2 = (x, z)$, $\mathfrak{m} = (x, y, z)$; $\mathfrak{p}_1$ and $\mathfrak{p}_2$ are prime, and $\mathfrak{m}$ is maximal. Let $\mathfrak{a} = \mathfrak{p}_1 \mathfrak{p}_2$. Show that $\mathfrak{a} = \mathfrak{p}_1 \cap \mathfrak{p}_2 \cap \mathfrak{m}^2$ is a reduced primary decomposition of $\mathfrak{a}$. Which components are isolated and which are embedded?

    \item \label{exer6-chapter4}Let $X$ be an infinite compact Hausdorff space, $C(X)$ the ring of real-valued continuous functions on $X$ (Chapter \ref{chapter1}, Exercise \ref{exer26-chapter1}). Is the zero ideal decomposable in this ring?

    \item \label{exer7-chapter4}Let $A$ be a ring and let $A[x]$ denote the ring of polynomials in one indeterminate over $A$. For each ideal $\mathfrak{a}$ of $A$, let $\mathfrak{a}[x]$ denote the set of all polynomials in $A[x]$ with coefficients in $\mathfrak{a}$.
    \begin{enumerate}[i)]
        \item $\mathfrak{a}[x]$ is the extension of $\mathfrak{a}$ to $A[x]$.
        \item If $\mathfrak{p}$ is a prime ideal in $A$, then $\mathfrak{p}[x]$ is a prime ideal in $A[x]$.
        \item If $\mathfrak{q}$ is a $\mathfrak{p}$-primary ideal in $A$, then $\mathfrak{q}[x]$ is a $\mathfrak{p}[x]$-primary ideal in $A[x]$.\\
        \textit{Hint:} Use Chapter 1, Exercise 2.
        \item If $\mathfrak{a} = \bigcap_{i=1}^n \mathfrak{q}_i$ is a minimal primary decomposition in $A$, then $\mathfrak{a}[x] = \bigcap_{i=1}^n \mathfrak{q}_i[x]$ is a minimal primary decomposition in $A[x]$.
        \item If $\mathfrak{p}$ is a minimal prime ideal of $\mathfrak{a}$, then $\mathfrak{p}[x]$ is a minimal prime ideal of $\mathfrak{a}[x]$.
    \end{enumerate}

    \item \label{exer8-chapter4}Let $k$ be a field. Show that in the polynomial ring $k[x_1, \ldots, x_n]$ the ideals $\mathfrak{p}_i = (x_1, \ldots, x_i)$ ($1 \leq i \leq n$) are prime and all their powers are primary.\\
    \textit{Hint:} Use Exercise \ref{exer7-chapter4}.

    \item \label{exer9-chapter4}In a ring $A$, let $D(A)$ denote the set of prime ideals $\mathfrak{p}$ which satisfy the following condition: there exists $a \in A$ such that $\mathfrak{p}$ is minimal in the set of prime ideals containing $(0 : a)$. Show that $x \in A$ is a zero divisor $\iff$ $x \in \mathfrak{p}$ for some $\mathfrak{p} \in D(A)$.\\
    Let $S$ be a multiplicatively closed subset of $A$, and identify $\Spec(S^{-1}A)$ with its image in $\Spec(A)$ (Chapter \ref{chapter3}, Exercise \ref{exer21-chapter3}). Show that
    \[
    D(S^{-1}A) = D(A) \cap \Spec(S^{-1}A).
    \]
    If the zero ideal has a primary decomposition, show that $D(A)$ is the set of associated prime ideals of $0$.

    \item \label{exer10-chapter4}For any prime ideal $\mathfrak{p}$ in a ring $A$, let $S_{\mathfrak{p}}(0)$ denote the kernel of the homomorphism $A \to A_{\mathfrak{p}}$. Prove that:
    \begin{enumerate}[i)]
        \item $S_{\mathfrak{p}}(0) \subseteq \mathfrak{p}$.
        \item $r(S_{\mathfrak{p}}(0)) = \mathfrak{p}$ $\iff$ $\mathfrak{p}$ is a minimal prime ideal of $A$.
        \item If $\mathfrak{p} \supseteq \mathfrak{p}'$, then $S_{\mathfrak{p}}(0) \subseteq S_{\mathfrak{p}'}(0)$.
        \item $\bigcap_{\mathfrak{p} \in D(A)} S_{\mathfrak{p}}(0) = 0$, where $D(A)$ is defined in Exercise 9.
    \end{enumerate}

    \item \label{exer11-chapter4}If $\mathfrak{p}$ is a minimal prime ideal of a ring $A$, show that $S_{\mathfrak{p}}(0)$ (Exercise \ref{exer10-chapter4}) is the smallest $\mathfrak{p}$-primary ideal.\\
    Let $\mathfrak{a}$ be the intersection of the ideals $S_{\mathfrak{p}}(0)$ as $\mathfrak{p}$ runs through the minimal prime ideals of $A$. Show that $\mathfrak{a}$ is contained in the nilradical of $A$.\\
    Suppose that the zero ideal is decomposable. Prove that $\mathfrak{a} = 0$ if and only if every prime ideal of $0$ is isolated.

    \item \label{exer12-chapter4}Let $A$ be a ring, $S$ a multiplicatively closed subset of $A$. For any ideal $\mathfrak{a}$, let $S(\mathfrak{a})$ denote the contraction of $S^{-1}\mathfrak{a}$ in $A$. The ideal $S(\mathfrak{a})$ is called the \emph{saturation}\index{saturation of ideal} of $\mathfrak{a}$ with respect to $S$. Prove that:
    \begin{enumerate}[i)]
        \item $S(\mathfrak{a}) \cap S(\mathfrak{b}) = S(\mathfrak{a} \cap \mathfrak{b})$.
        \item $S(r(\mathfrak{a})) = r(S(\mathfrak{a}))$.
        \item $S(\mathfrak{a}) = (1)$ $\iff$ $\mathfrak{a}$ meets $S$.
        \item $S_1(S_2(\mathfrak{a})) = (S_1 S_2)(\mathfrak{a})$.
    \end{enumerate}
    If $\mathfrak{a}$ has a primary decomposition, prove that the set of ideals $S(\mathfrak{a})$ (where $S$ runs through all multiplicatively closed subsets of $A$) is finite.

    \item \label{exer13-chapter4}Let $A$ be a ring and $\mathfrak{p}$ a prime ideal of $A$. The $n$th \emph{symbolic power}\index{symbolic power} of $\mathfrak{p}$ is defined to be the ideal (in the notation of Exercise \ref{exer12-chapter4})
    \[
    \mathfrak{p}^{(n)} = S_{\mathfrak{p}}(\mathfrak{p}^n)
    \]
    where $S_{\mathfrak{p}} = A \setminus \mathfrak{p}$. Show that:
    \begin{enumerate}[i)]
        \item $\mathfrak{p}^{(n)}$ is a $\mathfrak{p}$-primary ideal;
        \item if $\mathfrak{p}^n$ has a primary decomposition, then $\mathfrak{p}^{(n)}$ is its $\mathfrak{p}$-primary component;
        \item if $\mathfrak{p}^{(m)} \mathfrak{p}^{(n)}$ has a primary decomposition, then $\mathfrak{p}^{(m+n)}$ is its $\mathfrak{p}$-primary component;
        \item $\mathfrak{p}^{(n)} = \mathfrak{p}^n$ $\iff$ $\mathfrak{p}^{(n)}$ is $\mathfrak{p}$-primary.
    \end{enumerate}

    \item \label{exer14-chapter4}Let $\mathfrak{a}$ be a decomposable ideal in a ring $A$ and let $\mathfrak{p}$ be a maximal element of the set of ideals $(\mathfrak{a} : x)$, where $x \in A$ and $x \notin \mathfrak{a}$. Show that $\mathfrak{p}$ is a prime ideal belonging to $\mathfrak{a}$.

    \item \label{exer15-chapter4}Let $\mathfrak{a}$ be a decomposable ideal in a ring $A$, let $\Sigma$ be an isolated set of prime ideals belonging to $\mathfrak{a}$, and let $\mathfrak{q}_{\Sigma}$ be the intersection of the corresponding primary components. Let $f$ be an element of $A$ such that, for each prime ideal $\mathfrak{p}$ belonging to $\mathfrak{a}$, we have $f \in \mathfrak{p}$ $\iff$ $\mathfrak{p} \notin \Sigma$, and let $S_f$ be the set of all powers of $f$. Show that $\mathfrak{q}_{\Sigma} = S_f(\mathfrak{a}) = (\mathfrak{a} : f^n)$ for all large $n$.

    \item \label{exer16-chapter4}If $A$ is a ring in which every ideal has a primary decomposition, show that every ring of fractions $S^{-1}A$ has the same property.

    \item \label{exer17-chapter4}Let $A$ be a ring with the following property:
    \begin{enumerate}
        \item[(L1)] For every ideal $\mathfrak{a} \neq (1)$ in $A$ and every prime ideal $\mathfrak{p}$, there exists $x \notin \mathfrak{p}$ such that $S_{\mathfrak{p}}(\mathfrak{a}) = (\mathfrak{a} : x)$, where $S_{\mathfrak{p}} = A \setminus \mathfrak{p}$.
    \end{enumerate}
    \quad\, Then every ideal in $A$ is an intersection of (possibly infinitely many) primary ideals.

    \textit{Hint:} Let $\mathfrak{a}$ be an ideal $\neq (1)$ in $A$, and let $\mathfrak{p}_1$ be a minimal element of the set of prime ideals containing $\mathfrak{a}$. Then $\mathfrak{q}_1 = S_{\mathfrak{p}_1}(\mathfrak{a})$ is $\mathfrak{p}_1$-primary (by Exercise \ref{exer11-chapter4}), and $\mathfrak{q}_1 = (\mathfrak{a} : x)$ for some $x \notin \mathfrak{p}_1$. Show that $\mathfrak{a} = \mathfrak{q}_1 \cap (\mathfrak{a} + (x))$. 
    
    \quad\, Now let $\mathfrak{a}_1$ be a maximal element of the set of ideals $\mathfrak{b} \supseteq \mathfrak{a}$ such that $\mathfrak{q}_1 \cap \mathfrak{b} = \mathfrak{a}$, and choose $\mathfrak{a}_1$ so that $x \in \mathfrak{a}_1$, and therefore $\mathfrak{a}_1 \not\subseteq \mathfrak{p}_1$. Repeat the construction starting with $\mathfrak{a}_1$, and so on. At the $n$th stage we have $\mathfrak{a} = \mathfrak{q}_1 \cap \cdots \cap \mathfrak{q}_n \cap \mathfrak{a}_n$ where the $\mathfrak{q}_i$ are primary ideals, $\mathfrak{a}_n$ is maximal among the ideals $\mathfrak{b}$ containing $\mathfrak{a}_{n-1} = \mathfrak{a}_n \cap \mathfrak{q}_n$ such that $\mathfrak{a} = \mathfrak{q}_1 \cap \cdots \cap \mathfrak{q}_n \cap \mathfrak{b}$, and $\mathfrak{a}_n \not\subseteq \mathfrak{p}_n$. If at any stage we have $\mathfrak{a}_n = (1)$, the process stops, and $\mathfrak{a}$ is a finite intersection of primary ideals. If not, continue by transfinite induction, observing that each $\mathfrak{a}_n$ strictly contains $\mathfrak{a}_{n-1}$.

    \item \label{exer18-chapter4}Consider the following condition on a ring $A$:
    \begin{enumerate}
        \item[(L2)] Given an ideal $\mathfrak{a}$ and a descending chain $S_1 \supseteq S_2 \supseteq \cdots \supseteq S_n \supseteq \cdots$ of multiplicatively closed subsets of $A$, there exists an integer $n$ such that $S_n(\mathfrak{a}) = S_{n+1}(\mathfrak{a}) = \cdots$.
    \end{enumerate}
    Prove that the following are equivalent:
    \begin{enumerate}[i)]
        \item Every ideal in $A$ has a primary decomposition;
        \item $A$ satisfies (L1) and (L2).
    \end{enumerate}
    \textit{Hint:} For i) $\implies$ ii), use Exercises \ref{exer12-chapter4} and \ref{exer15-chapter4}. For ii) $\implies$ i) show, with the notation of the proof of Exercise \ref{exer17-chapter4}, that if $S_n = S_{\mathfrak{p}_1} \cap \cdots \cap S_{\mathfrak{p}_n}$ then $S_n$ meets $\mathfrak{a}_n$, hence $S_n(\mathfrak{a}_n) = (1)$, and therefore $S_n(\mathfrak{a}) = \mathfrak{q}_1 \cap \cdots \cap \mathfrak{q}_n$. Now use (L2) to show that the construction must terminate after a finite number of steps.

    \item \label{exer19-chapter4}Let $A$ be a ring and $\mathfrak{p}$ a prime ideal of $A$. Show that every $\mathfrak{p}$-primary ideal contains $S_{\mathfrak{p}}(0)$, the kernel of the canonical homomorphism $A \to A_{\mathfrak{p}}$.
    
    \quad\, Suppose that $A$ satisfies the following condition: for every prime ideal $\mathfrak{p}$, the intersection of all $\mathfrak{p}$-primary ideals of $A$ is equal to $S_{\mathfrak{p}}(0)$. (Noetherian rings satisfy this condition: see Chapter \ref{chapter10}.) Let $\mathfrak{p}_1, \ldots, \mathfrak{p}_n$ be distinct prime ideals, none of which is a minimal prime ideal of $A$. Then there exists an ideal $\mathfrak{a}$ in $A$ whose associated prime ideals are $\mathfrak{p}_1, \ldots, \mathfrak{p}_n$.

    \textit{Hint:} Proof by induction on $n$. The case $n = 1$ is trivial (take $\mathfrak{a} = \mathfrak{p}_1$). Suppose $n > 1$ and let $\mathfrak{p}_n$ be maximal in the set $\{\mathfrak{p}_1, \ldots, \mathfrak{p}_n\}$. By the inductive hypothesis there exists an ideal $\mathfrak{b}$ and a minimal primary decomposition $\mathfrak{b} = \mathfrak{q}_1 \cap \cdots \cap \mathfrak{q}_{n-1}$, where each $\mathfrak{q}_i$ is $\mathfrak{p}_i$-primary. If $\mathfrak{b} \not\subseteq S_{\mathfrak{p}_n}(0)$, let $\mathfrak{p}$ be a minimal prime ideal of $A$ contained in $\mathfrak{p}_n$. Then $S_{\mathfrak{p}_n}(0) \not\subseteq S_{\mathfrak{p}}(0)$, hence $\mathfrak{b} \not\subseteq S_{\mathfrak{p}}(0)$. Taking radicals and using Exercise \ref{exer10-chapter4}, we have $\mathfrak{p}_1 \cap \cdots \cap \mathfrak{p}_{n-1} \subseteq \mathfrak{p}$, hence some $\mathfrak{p}_i \subseteq \mathfrak{p}$, hence $\mathfrak{p}_i = \mathfrak{p}$ since $\mathfrak{p}$ is minimal. This is a contradiction since no $\mathfrak{p}_i$ is minimal. Hence $\mathfrak{b} \not\subseteq S_{\mathfrak{p}_n}(0)$ and therefore there exists a $\mathfrak{p}_n$-primary ideal $\mathfrak{q}_n$ such that $\mathfrak{b} \not\subseteq \mathfrak{q}_n$. Show that $\mathfrak{a} = \mathfrak{q}_1 \cap \cdots \cap \mathfrak{q}_n$ has the required properties.
\end{enumerate}

\emph{Primary decomposition of modules}

Practically the whole of this chapter can be transposed to the context of modules over a ring $A$. The following exercises indicate how this is done.
\begin{enumerate}
    \setcounter{enumi}{19}
    \item \label{exer20-chapter4}Let $M$ be a fixed $A$-module, $N$ a submodule of $M$. The \emph{radical of $N$ in $M$}\index{radical!of submodule} is defined to be
    \[
    r_M(N) = \{x \in A : x^q M \subseteq N \text{ for some } q > 0\}.
    \]
    Show that $r_M(N) = r(N : M) = r(\Ann(M/N))$. In particular, $r_M(N)$ is an ideal.
    
    \quad\, State and prove the formulas for $r_M$ analogous to \eqref{exer:radical-properties}.

    \item \label{exer21-chapter4}An element $x \in A$ defines an endomorphism $\phi_x$ of $M$, namely $m \mapsto xm$. The element $x$ is said to be a \emph{zero-divisor} (resp. \emph{nilpotent}) in $M$ if $\phi_x$ is not injective (resp. is nilpotent). A submodule $Q$ of $M$ is \emph{primary in $M$}\index{primary submodule} if $Q \neq M$ and every zero-divisor in $M/Q$ is nilpotent.
    
    \quad\, Show that if $Q$ is primary in $M$, then $(Q : M)$ is a primary ideal and hence $r_M(Q)$ is a prime ideal. We say that $Q$ is $\mathfrak{p}$-primary (in $M$).

    \quad\, Prove the analogues of \eqref{lem:intersection-primary} and \eqref{lem:colon-primary}.

    \item \label{exer22-chapter4}A \emph{primary decomposition of $N$ in $M$} is a representation of $N$ as an intersection
    \[
    N = Q_1 \cap \cdots \cap Q_n
    \]
    of primary submodules of $M$; it is a minimal primary decomposition if the ideals $\mathfrak{p}_i = r_M(Q_i)$ are all distinct and if none of the components $Q_i$ can be omitted from the intersection, that is if $Q_i \not\supseteq \bigcap_{j \neq i} Q_j$ ($1 \leq i \leq n$).

    Prove the analogue of \eqref{thm:first-uniqueness}, that the prime ideals $\mathfrak{p}_i$ depend only on $N$ (and $M$). They are called the prime ideals belonging to $N$ in $M$. Show that they are also the prime ideals belonging to $0$ in $M/N$.

    \item \label{exer23-chapter4}State and prove the analogues of \eqref{prop:minimal-primes}--\eqref{cor:isolated-primary-components-unique} inclusive. (There is no loss of generality in taking $N = 0$.)
\end{enumerate}

